% Owner: Kashaf Ali / Member A. Feedback revision approved 10 October 2026.
Zeest addresses the difficulty of reviewing patient information scattered across consultations, medical images, laboratory investigations and medication records. Physicians need relevant history alongside the current encounter, but preparing that context and documenting the consultation require additional effort. The intended application brings these tasks into one web workspace. Each physician accesses their own patient records, selects an encounter and reviews its historical context. Consultation audio becomes structured notes that can be checked and corrected before analysis. Supported imaging, laboratory information and medication context contribute identifiable findings rather than independent clinical decisions. Longitudinal memory combines explicit patient and encounter relationships with semantic retrieval, preserving the source of information across visits. Specialised agents organise the case, process supported inputs and retrieve relevant clinical guidance. The resulting diagnosis and treatment draft appears with cited evidence and critic findings so the physician can inspect its support and limitations. This workflow aims to make information continuity and recommendation review easier to follow while keeping clinical judgement with the physician. Its value requires assessment through implementation and controlled evaluation rather than assumptions about automation.

Approval, modification and rejection have different consequences in Zeest. Approval records the exact reviewed revision as a clinical decision; modification saves a pending revision, and rejection retains the review action without accepting the draft. Audit history preserves the physician, encounter, time and reviewed content, while patient ownership governs related information access. The current deliverable establishes the project vision, requirements, use cases, interface designs, database relationships, methodology and software architecture. It does not claim an implemented clinical system or measured improvements in patient outcomes. Development and evaluation require permitted de-identified or synthetic cases, suitable existing models and selected knowledge resources. Diagnostic models will be integrated rather than trained from scratch, and clinical coverage will follow their supported tasks. The evaluation method compares a single language model, a memory-augmented variant and the full system, with controlled contrasts for critic and multimodal contributions. These comparisons will determine which components provide useful support under common assessment conditions. Autonomous treatment, real clinical deployment and regulatory approval remain outside scope; every accepted decision requires physician oversight.
